% ======================================================================
\section{Details of the single-gluon incoming state}\label{app:onegluon}
% ======================================================================

This appendix contains the intermediate steps of Section~\ref{sec:onegin}: the Fourier transforms, the explicit forms of the terms with the three-gluon vertex, the cancellation of the double pole, the algebra of the terms with two valence vertices, the coordinate-space forms of all components, and the verification of the normalisation. The organisation follows that of the supplementary Appendix~G of~\cite{Lublinsky:2016meo} for the vacuum initial state.

\subsection{Fourier transforms}\label{app:fourier}

The Fourier transforms that are needed are, for real $\xi>0$ and real $a,b$,
\begin{align}
&\int\frac{d^2k}{2\pi}\,\frac{k^i}{\v k^2}\,e^{i\v k\cdot\v X}=\frac{iX^i}{\v X^2}\,,\label{eq:FT_K}\\
&\int\frac{d^2k}{2\pi}\frac{d^2p}{2\pi}\,\frac{k^ip^j}{\v k^2\,(a\v k^2+b\v p^2)}\,e^{i\v k\cdot\v X+i\v p\cdot\v Z}=-\frac{X^iZ^j}{\v Z^2\,(b\v X^2+a\v Z^2)}\,,\label{eq:FT_i2}\\
&\int\frac{d^2k}{2\pi}\frac{d^2p}{2\pi}\,\frac{1}{\xi\v k^2-\v p^2}\,e^{i\v k\cdot\v X+i\v p\cdot\v Z}=\frac{1}{\xi\v Z^2-\v X^2}\,,\label{eq:FT_i3}\\
&\int\frac{d^2k}{2\pi}\frac{d^2p}{2\pi}\,\frac{k^m-p^m}{(\v k-\v p)^2}\,\frac{p^l-\xi k^l}{\v p^2-\xi\v k^2}\,e^{i\v k\cdot\v X+i\v p\cdot\v Z}=-\frac{X^m+\xi Z^m}{\v X^2-\xi\v Z^2}\,\frac{X^l+Z^l}{(\v X+\v Z)^2}\,,\label{eq:FT_i4}\\
&\int\frac{d^2k}{2\pi}\frac{d^2p}{2\pi}\,\frac{k^m-p^m}{\v p^2-\xi\v k^2}\,\frac{p^l-\xi k^l}{(\v p-\xi\v k)^2}\,e^{i\v k\cdot\v X+i\v p\cdot\v Z}=-\frac{X^l+Z^l}{\v X^2-\xi\v Z^2}\,\frac{X^m+\xi Z^m}{(\v X+\xi\v Z)^2}\,,\label{eq:FT_i5}
\end{align}
with principal values where a denominator vanishes~\cite{Lublinsky:2016meo,Radhakrishnan:2026soft}. In all cases below they are applied with $\v X\to\v X'$ and $\v Z\to-\v X$ or $\v Z\to-\v Y$, because the incoming gluon contributes the phase $e^{-i\v k\cdot\v z'}$, the outgoing gluon the phase $e^{i\v p\cdot\v z}$ and a charge $\rho(\v q)=\int_xe^{i\v q\cdot\v x}\rho(\v x)$.

Eq.~\eqref{eq:FT_K} follows from $\int_0^{2\pi}d\theta\,\cos\theta\,e^{ikX\cos\theta}=2\pi iJ_1(kX)$ and $\int_0^\infty dt\,J_1(t)=1$. Eqs.~\eqref{eq:FT_i2} and~\eqref{eq:FT_i3} follow from Eq.~\eqref{eq:FT_K} and from the transform of $1/(a\v k^2+b\v p^2)$ in four dimensions.

\NOTE{Eqs.~\eqref{eq:FT_i4} and~\eqref{eq:FT_i5} are quoted from the sources and were not re-derived here.}

\subsection{Vacuum outgoing state}

For Eq.~\eqref{eq:psi_rho}: $\braket{0|H_g|g^a_i(k)}=g\,k^i\rho^a(\v k)/[4\pi^{3/2}(k^+)^{3/2}]$ and $E_g(k)=\v k^2/2k^+$. In coordinate space, with Eq.~\eqref{eq:FT_K},
\begin{align}\label{eq:psi_rhoX}
\ket{(\Psi^a_i(k^+,\v z'))_\rho}=\frac{ig}{2\pi^{3/2}\sqrt{k^+}}\int_x\frac{X'^i}{\v X'^2}\,\rho^a(\v x)\ket0\,.
\end{align}

\subsection{One outgoing gluon: instantaneous term}

The matrix element in Eq.~\eqref{eq:psi_inst_def} is Eq.~\eqref{eq:m5} with $(p,b,i)\to(k,a,i)$ and $(q,c,j)\to(p,b,j)$, i.e.\ $ig^2f^{abc}(k^++p^+)\delta_{ij}\rho^c(\v k-\v p)/[2(2\pi)^3\sqrt{k^+p^+}(k^+-p^+)^2]$ after renaming the colour index of the charge. Dividing by $E_g(k)-E_g(p)$ gives Eq.~\eqref{eq:psi_inst}. With the momentum fraction $\xi\equiv p^+/k^+$ and Eq.~\eqref{eq:FT_i3},
\begin{align}\label{eq:psi_instX}
\ket{(\Psi^a_i(k^+,\v z'))_{g\rho\,{\rm inst}}}=\frac{ig^2f^{abc}}{(2\pi)^3}\int d\xi\int_{x,z}\frac{\sqrt\xi\,(1+\xi)}{(1-\xi)^2}\,\frac{\rho^c(\v x)}{\xi\v X^2-\v X'^2}\,\ket{g^b_i(\xi k^+,\v z)}\,,
\end{align}
where the $\xi$ integral runs over $1+\Lambda/k^+<\xi<\vee/k^+$ in region~1 and over $\Lambda/k^+<\xi<1-\Lambda/k^+$ in region~2. Note that $\xi>1$ in region~1.


\subsection{One outgoing gluon: terms with the three-gluon vertex}

In Eq.~\eqref{eq:psi1_def} the matrix element of $H_g$, Eq.~\eqref{eq:m2}, sets one of the intermediate gluons equal to the incoming one, $q=k$, $(c,l)=(a,i)$, and the other to the emitted gluon, $r=p-k$; the conjugate of Eq.~\eqref{eq:m3} with the parent $(p,b,j)$ is $-igf^{bda}\,[\dots]/[8\pi^{3/2}\sqrt{p^+r^+k^+}]$, and the denominators are $-E_g(r)$ and $E_g(k)-E_g(p)$. Inserting the matrix elements we find
\begin{align}\label{eq:psi1}
\ket{(\Psi^a_i(k))^1_{g\rho}}=\ &\frac{ig^2f^{abc}}{8\pi^3}\int_{k^++\Lambda}^{\vee}dp^+\!\int d^2p\ \frac{\rho^c(\v k-\v p)\,(p-k)^l}{\sqrt{p^+k^+}\,(p^+-k^+)^2}\,
\frac{1}{\Big(\frac{\v k^2}{k^+}-\frac{\v p^2}{p^+}\Big)\frac{(\v k-\v p)^2}{k^+-p^+}}\notag\\
&\times\Bigg[\Big(2k^j-p^j+\frac{p^+-2k^+}{p^+}\,p^j\Big)\delta_{il}+\Big(2p^i-k^i-\frac{2p^+-k^+}{k^+}\,k^i\Big)\delta_{jl}\notag\\
&\qquad+\Big(\frac{p^++k^+}{p^+-k^+}\,(p-k)^l-k^l-p^l\Big)\delta_{ij}\Bigg]\ket{g^b_j(p)}\,.
\end{align}
The three brackets are $-\frac2\xi(p-\xi k)^j$, $2(p-\xi k)^i$ and $-\frac{2}{1-\xi}(p-\xi k)^l$ with $\xi=p^+/k^+$, so that
Eq.~\eqref{eq:psi1_xi} follows. With Eq.~\eqref{eq:FT_i4},
\begin{align}\label{eq:psi1X}
\ket{(\Psi^a_i(k^+,\v z'))^1_{g\rho}}=\ &-\frac{ig^2f^{abc}}{4\pi^3}\int_{1+\Lambda/k^+}^{\vee/k^+}d\xi\int_{x,z}\frac{\sqrt\xi}{1-\xi}\,\frac{X'^m-\xi X^m}{\v X'^2-\xi\v X^2}\,\frac{X'^l-X^l}{(\v X'-\v X)^2}\,\rho^c(\v x)\notag\\
&\times\Big(\delta_{il}\delta_{jm}-\frac1\xi\,\delta_{jl}\delta_{im}-\frac{1}{1-\xi}\,\delta_{ij}\delta_{lm}\Big)\ket{g^b_j(\xi k^+,\v z)}\,.
\end{align}
In Eq.~\eqref{eq:psi2_def} the matrix element of $H_{ggg}$ is Eq.~\eqref{eq:m3} with the outgoing gluons $p$ and $r=k-p$, and the matrix element of $H_g$ is the conjugate of Eq.~\eqref{eq:m2}, which absorbs the gluon $r$ with the factor $g\,r^m\rho^d(\v r)/[4\pi^{3/2}(r^+)^{3/2}]$. The result is
\begin{align}\label{eq:psi2}
\ket{(\Psi^a_i(k))^2_{g\rho}}=\ &\frac{ig^2f^{abc}}{8\pi^3}\int_{\Lambda}^{k^+-\Lambda}dp^+\!\int d^2p\ \frac{\rho^c(\v k-\v p)\,(k-p)^l}{\sqrt{k^+p^+}\,(k^+-p^+)^2}\notag\\
&\times\frac{1}{\Big(\frac{\v k^2}{k^+}-\frac{(\v k-\v p)^2}{k^+-p^+}-\frac{\v p^2}{p^+}\Big)\Big(\frac{\v k^2}{k^+}-\frac{\v p^2}{p^+}\Big)}\notag\\
&\times\Bigg[\Big(2p^i-k^i+\frac{k^+-2p^+}{k^+}\,k^i\Big)\delta_{jl}+\Big(2k^j-p^j-\frac{2k^+-p^+}{p^+}\,p^j\Big)\delta_{il}\notag\\
&\qquad+\Big(\frac{k^++p^+}{k^+-p^+}\,(k-p)^l-p^l-k^l\Big)\delta_{ij}\Bigg]\ket{g^b_j(p)}\,,
\end{align}
or, with $\frac{\v k^2}{k^+}-\frac{(\v k-\v p)^2}{k^+-p^+}-\frac{\v p^2}{p^+}=-\frac{(\v p-\xi\v k)^2}{\xi(1-\xi)k^+}$,
Eq.~\eqref{eq:psi2_xi} follows. In coordinate space, with Eq.~\eqref{eq:FT_i5},
\begin{align}\label{eq:psi2X}
\ket{(\Psi^a_i(k^+,\v z'))^2_{g\rho}}=\ &-\frac{ig^2f^{abc}}{4\pi^3}\int_{\Lambda/k^+}^{1-\Lambda/k^+}d\xi\ \frac{\xi^{3/2}}{1-\xi}\int_{x,z}\frac{X'^l-X^l}{\v X'^2-\xi\v X^2}\,\frac{X'^m-\xi X^m}{(\v X'-\xi\v X)^2}\,\rho^c(\v x)\notag\\
&\times\Big(\delta_{il}\delta_{jm}-\frac1\xi\,\delta_{jl}\delta_{im}-\frac{1}{1-\xi}\,\delta_{lm}\delta_{ij}\Big)\ket{g^b_j(\xi k^+,\v z)}\,.
\end{align}

\subsection{Cancellation of the double pole}\label{app:doublepole}

Each of the four terms linear in $\rho$ has a double pole at $p^+=k^+$, i.e.\ a factor $1/(1-\xi)^2$. The pole lies outside the regions~\eqref{eq:ranges}, but the limits of integration approach it from both sides at the distance $\Lambda$, and a single term gives a contribution proportional to $k^+/\Lambda$. In the sum of the instantaneous and the non-instantaneous terms the double pole cancels in each region separately: for $p^+\to k^+$ the last bracket of Eq.~\eqref{eq:psi1} is dominated by its third term, and one finds that the integrand of Eq.~\eqref{eq:psi1} approaches minus the integrand of Eq.~\eqref{eq:psi_inst}; the same holds for Eq.~\eqref{eq:psi2}. No term proportional to $k^+/\Lambda$ remains.

Numerically, for a generic configuration of the transverse momenta, $(p^+-k^+)^2$ times the sum of the instantaneous and the non-instantaneous term of Eq.~\eqref{eq:B3mom} is $1.7\times10^{-3}$, $1.7\times10^{-4}$, $1.7\times10^{-5}$ for $|p^+-k^+|/k^+=10^{-2},10^{-3},10^{-4}$ on either side, while the instantaneous term alone gives $1.46$ (script \texttt{scripts/check\_onegluon\_state.py}). A simple pole remains; it is odd in $p^+-k^+$ and gives no term $\propto k^+/\Lambda$.

\subsection{Terms with two valence vertices}\label{app:tworho}

In Eq.~\eqref{eq:psi1rr} the matrix elements are the conjugate of Eq.~\eqref{eq:m1} for the absorption and Eq.~\eqref{eq:m1} for the emission, and $E_g(k)\,(E_g(k)-E_g(p))=\v k^2(p^+\v k^2-k^+\v p^2)/(4k^{+2}p^+)$. In Eq.~\eqref{eq:psi2rr} the two terms of Eq.~\eqref{eq:m2} and of its conjugate give two equal connected contributions, which cancel the factor $\frac12$, and the disconnected one discussed in the main text. The sum of the two contributions is
\begin{align}
\frac{g^2\,k^ip^j}{4\pi^3\sqrt{k^+p^+}\,\v k^2\v p^2}\ \frac{k^+\v p^2\,\rho^b(-\v p)\rho^a(\v k)-p^+\v k^2\,\rho^a(\v k)\rho^b(-\v p)}{p^+\v k^2-k^+\v p^2}\,.
\end{align}
With the decomposition~\eqref{eq:decomp}, $\rho^a(\v k)\rho^b(-\v p)=\frac12\{\rho^a(\v k),\rho^b(-\v p)\}+\frac i2f^{abc}\rho^c(\v k-\v p)$ and the opposite sign of the last term for the opposite order, the energy denominator cancels in the anticommutator part, which gives Eqs.~\eqref{eq:rrsum}--\eqref{eq:psi_comm}.
In coordinate space, with Eqs.~\eqref{eq:FT_K} and~\eqref{eq:FT_i2} and $\xi=p^+/k^+$,
\begin{align}
\ket{(\Psi^a_i(k^+,\v z'))_{g\{\rho,\rho\}}}&=-\frac{g^2}{8\pi^3}\int_{\Lambda/k^+}^{\vee/k^+}\frac{d\xi}{\sqrt\xi}\int_{x,y,z}\frac{X'^i\,Y^j}{\v X'^2\,\v Y^2}\,\big\{\rho^a(\v x),\rho^b(\v y)\big\}\ket{g^b_j(\xi k^+,\v z)}\,,\label{eq:psi_antiX}\\
\ket{(\Psi^a_i(k^+,\v z'))_{g[\rho,\rho]}}&=-\frac{ig^2f^{abc}}{8\pi^3}\int_{\Lambda/k^+}^{\vee/k^+}\frac{d\xi}{\sqrt\xi}\int_{x,z}\frac{X'^i\,X^j}{\v X'^2\,\v X^2}\,\frac{\v X'^2+\xi\v X^2}{\xi\v X^2-\v X'^2}\,\rho^c(\v x)\ket{g^b_j(\xi k^+,\v z)}\,.\label{eq:psi_commX}
\end{align}


\subsection{Two outgoing gluons}

In coordinate space
\begin{align}\label{eq:psiggrhoX}
\ket{(\Psi^a_i(k^+,\v z'))_{gg\rho}}=\frac{ig}{2\pi^{3/2}}\int_\Lambda^\vee\frac{dp^+}{\sqrt{p^+}}\int_{x,z}\frac{X^j}{\v X^2}\,\rho^b(\v x)\,\ket{g^a_i(k^+,\v z')\,g^b_j(p^+,\v z)}\,.
\end{align}
For the $\rho$-independent component~\eqref{eq:psigg} the coordinate-space form is contained in the coefficient $\mf C_1$, Eq.~\eqref{eq:C1X}.

\subsection{Normalisation of the dressed one-gluon state}\label{app:norm}

We verify Eq.~\eqref{eq:normcond} to order $g^2$. The terms of order $g^2$ in the overlap are the overlaps of two components of order $g$ and the overlaps of $\ket g$ with the components of order $g^2$.

\paragraph{Terms with two charges.} From Eqs.~\eqref{eq:psi_rho} and~\eqref{eq:psiggrho}, with the normalisation~\eqref{eq:norms} of the two-gluon states,
\begin{align}
\braket{(\Psi^{a'}_{i'}(k'))_\rho|(\Psi^a_i(k))_\rho}&=\frac{g^2}{4\pi^3\sqrt{k^+k'^+}}\,\frac{k^ik'^{i'}}{\v k^2\v k'^2}\,\rho^{a'}(-\v k')\,\rho^a(\v k)\,,\label{eq:ov1}\\
\braket{(\Psi^{a'}_{i'}(k'))_{gg\rho}|(\Psi^a_i(k))_{gg\rho}}&=\frac{g^2}{4\pi^3\sqrt{k^+k'^+}}\,\frac{k^ik'^{i'}}{\v k^2\v k'^2}\,\rho^a(\v k)\,\rho^{a'}(-\v k')\notag\\
&\quad+\delta^{aa'}\delta_{ii'}\delta^{(3)}(k-k')\,\frac{g^2}{4\pi^3}\int_\Lambda^\vee\frac{dp^+}{p^+}\int d^2p\,\frac{\rho^b(\v p)\,\rho^b(-\v p)}{\v p^2}\,,\label{eq:ov2}
\end{align}
where the first term of Eq.~\eqref{eq:ov2} comes from the exchange term of Eq.~\eqref{eq:norms}. From Eq.~\eqref{eq:psi_anti},
\begin{align}\label{eq:ov3}
\braket{g^{a'}_{i'}(k')|(\Psi^a_i(k))_{g\{\rho,\rho\}}}+\braket{(\Psi^{a'}_{i'}(k'))_{g\{\rho,\rho\}}|g^a_i(k)}=-\frac{g^2}{4\pi^3\sqrt{k^+k'^+}}\,\frac{k^ik'^{i'}}{\v k^2\v k'^2}\,\big\{\rho^a(\v k),\rho^{a'}(-\v k')\big\}\,.
\end{align}
The terms that are not proportional to $\delta^{(3)}(k-k')$ cancel in the sum of Eqs.~\eqref{eq:ov1}--\eqref{eq:ov3}, for non-commuting charges. The remaining term of Eq.~\eqref{eq:ov2} is $-2\mf N_2\,\delta^{aa'}\delta_{ii'}\delta^{(3)}(k-k')$ and is cancelled by the term $2\mf N_2$ from the first term of Eq.~\eqref{eq:psi_sum}, with $\mf N_2$ of Eq.~\eqref{eq:N2mom}. The coefficient of the identity in $\Omega$ is thus the same in the vacuum and in the one-gluon sector, as it must be.

\paragraph{Terms linear in $\rho$.} Consider a component of the form $\int d^3p\,\ket{g^b_j(p)}\,M^{ba}_{ji}(k,p)/(E_g(k)-E_g(p))$. Its overlap with $\bra{g^{a'}_{i'}(k')}$, plus the conjugate term, is
\begin{align}
\frac{M^{a'a}_{i'i}(k,k')-\big[M^{aa'}_{ii'}(k',k)\big]^\dagger}{E_g(k)-E_g(k')}\,.
\end{align}
It vanishes if $M$ is the matrix element of a Hermitian operator; this is the case for the instantaneous terms~\eqref{eq:psi_inst}. The commutator term~\eqref{eq:psi_comm} cancels against its conjugate in the same way, because its kernel is odd under the exchange of the two gluons. The terms with the three-gluon vertex, Eqs.~\eqref{eq:psi1} and~\eqref{eq:psi2}, do not cancel by themselves; they cancel against the overlaps of the two-gluon components $\ket{(\Psi)_{gg}}$ and $\ket{(\Psi)_{gg\rho}}$,
\begin{align}\label{eq:ov4}
&\braket{g^{a'}_{i'}(k')|(\Psi^a_i(k))^1_{g\rho}}+\braket{g^{a'}_{i'}(k')|(\Psi^a_i(k))^2_{g\rho}}\notag\\
&\quad+\braket{(\Psi^{a'}_{i'}(k'))^1_{g\rho}|g^a_i(k)}+\braket{(\Psi^{a'}_{i'}(k'))^2_{g\rho}|g^a_i(k)}\notag\\
&\quad+\braket{(\Psi^{a'}_{i'}(k'))_{gg\rho}|(\Psi^a_i(k))_{gg}}+\braket{(\Psi^{a'}_{i'}(k'))_{gg}|(\Psi^a_i(k))_{gg\rho}}=0\,.
\end{align}
Only the sum of the six terms vanishes. Eq.~\eqref{eq:ov4} is the triple-vertex part of the unitarity relation~\eqref{eq:B3unit} below, where it is derived; we have verified it numerically with Eqs.~\eqref{eq:B3mom}, \eqref{eq:C1mom} and~\eqref{eq:A1mom}.

\paragraph{Terms independent of $\rho$.} From Eq.~\eqref{eq:psigg}, summing over colours and polarisations,
\begin{align}\label{eq:ov5}
\braket{(\Psi^{a'}_{i'}(k'))_{gg}|(\Psi^a_i(k))_{gg}}&=\delta^{aa'}\delta_{ii'}\delta^{(3)}(k-k')\,\frac{g^2N_c}{8\pi^3}\int\frac{d^2\tilde p}{\tilde{\v p}^2}\notag\\
&\quad\times\int_{\Lambda/k^+}^{1-\Lambda/k^+}d\xi\,\Big[\xi(1-\xi)+\frac{1-\xi}{\xi}+\frac{\xi}{1-\xi}\Big]\notag\\
&=\delta^{aa'}\delta_{ii'}\delta^{(3)}(k-k')\,\frac{g^2N_c}{8\pi^3}\,\Big[2\log\frac{k^+}{\Lambda}-\frac{11}{6}\Big]\int\frac{d^2\tilde p}{\tilde{\v p}^2}\,,
\end{align}
with $\tilde{\v p}=\v p-\xi\v k$ and up to terms that vanish for $\Lambda/k^+\to0$. This is the probability for the incoming gluon to split. It has to be compensated by a $\rho$-independent term in the component along $\ket{g^a_i(k)}$. Such a term is not part of $\mf N$, which is fixed by the vacuum sector; it is a contribution to $\mf B_3$ proportional to $\delta^{(3)}(k-p)$,
\begin{align}\label{eq:B3rho0}
\Big[\mf B^{ba}_{3ji}(k,p)\Big]_{\rho^0}=-\frac12\,(2\pi)^3\,\delta^{ab}\delta_{ij}\,\delta^{(3)}(k-p)\,\frac{g^2N_c}{8\pi^3}\,\Big[2\log\frac{k^+}{\Lambda}-\frac{11}{6}\Big]\int\frac{d^2\tilde p}{\tilde{\v p}^2}\,,
\end{align}
which is the second term of Eq.~\eqref{eq:psi_sum} and the term $-\frac12\sum_mV_{im}V_{mi}/(E_i-E_m)^2$ of Eq.~\eqref{eq:pt2} with two-gluon intermediate states. It is the wave-function renormalisation of the soft gluon. The transverse integral $\int d^2\tilde p/\tilde{\v p}^2$ has no scale. It diverges logarithmically at small and at large $\tilde{\v p}$, and it vanishes in dimensional regularisation, where the two divergences cancel each other. With this rule Eq.~\eqref{eq:B3rho0} is zero, the $\rho$-independent term is absent from $\mf B_3$, Eq.~\eqref{eq:B3mom}, and
\begin{align}\label{eq:Nresult}
\mf N=1-\frac{g^2}{8\pi^3}\,\log\frac{\vee}{\Lambda}\int d^2p\,\frac{\rho^a(-\v p)\,\rho^a(\v p)}{\v p^2}
\end{align}
is the complete coefficient of $\ket{g^a_i(k)}$ at order $g^2$.
